\section{2024-12-16}

\subsection{Jean-Jacques Rousseau: Mensch, die Übel rühren von dir her!}

\startitemize[packed]
\item Der Mensch soll nicht nach übel suchen, er ist es selber
\item alles schlechte in der Welt kommt vom Menscheno
\item Gerechtigkeit und güte ist das unzertrennlcih
\item für den Menschen spielt nichts eine Rolle, ausser der eigenen Vorteil (Z. 19)
\item Gründung
\startitemize[packed]
\item Erster Mensch der etwas als \quotation{seins} annerkannte
\stopitemize
\stopitemize

\startplacetable[location={here},title={}]
\setupTABLE[r][1][background=color,backgroundcolor=gray]
\bTABLE[offset=1mm]
\startCSV
Darstellung der Welt, Darstellung des Menschen
\stopCSV
\bTR
\bTD % Darstellung Welt
\startitemize[packed]
\item Gut, ohne Übel
\item System der Ordnung -> Übel stammt aus Unordnung
\item Gerehtigkeit unzertrennlich mit Güte verbunden
\item Übel wird durch Menschen in die Welt getragen
\item Erde gehört niemanden, Früchte allen
\item Besitzt und Egoismus als Quelle als Ursprung allen Übel
\stopitemize
\eTD
\bTD % Dartellung Mensch
\startitemize[packed]
\item Einzieger Urheber der Übels
\item Natur schreibt ihm einfach, glecihförmige, solitäre Lebensweise vor
\item Im Herzen Moralisch
\startitemize[packed]
\item Strebt nach Liebe, Freundschaft, heroischen Taten
\stopitemize
\item Schafft ein system der Unordnung
\stopitemize
\eTD
\eTR
\eTABLE
\stopplacetable

\subsection{Die Natur des Bösen}

\startitemize[packed]
\item Das Böre fasziniert
\item Das Gute ist langweilig
\item Natur kennt keine Moral
\item Der Mensch beurteilt diese allerdings nach seinen eigenen Wertevorstellungen
\item Wettbewerb ist in der Natur allgegenwärtig
\item Aggressives Verhalten gegenüber Artgenossen normal
\item 
\stopitemize
